% ===== §4 Why Field-Building Is Not an Engineering Problem =====
\section{The Construction of a Field Is Not an Engineering Problem}
\label{p22:sec:not-engineering}

\noindent
This is the hardest link in the chain and the one on which the move to aesthetics depends. The field has been defined as a conditional environment of coupling, and the claim now to be argued is that the construction of such a field is an aesthetic and not an engineering matter. The claim is contested because it appears to deny something obvious, that conditions can be designed, that a person can arrange rhythm and atmosphere and the placement of things deliberately and skilfully, which sounds like exactly the sort of optimisation engineering performs. The section's task is to show why this appearance misleads, what it is about the construction of a field that escapes the engineering it superficially resembles, and it will do so by isolating the feature of field-building that no engineering can capture, the state-dependent appropriateness that cannot be standardised, and then by meeting the strongest objection, the one that points to architecture and design as disciplines that are at once engineering and aesthetic.

\subsection{The nature and the requirements of engineering}
\label{p22:sub:eng-what}

To argue that field-building is not engineering, one must say what engineering is, sharply enough that the contrast is real and not a caricature. Engineering, in the sense relevant here, is the discipline that solves problems by means that are repeatable, standardisable, and separable into ends and means. Its characteristic achievement is the specification: a procedure, a design, a set of parameters that, applied to a recurring type of situation, reliably produces a desired type of outcome, and that can be transmitted, taught, and applied by another without the transmitter present. An engineering solution to a problem of rhythm or arrangement would be a specification of the right rhythm, the right arrangement, the right distribution of attention, such that following the specification would build the field, as following a blueprint builds the bridge. The power of engineering lies exactly in this separability and repeatability: the bridge does not need the engineer to have a relationship with the river, and the specification works the same for every river of the relevant type. Engineering succeeds where the appropriate is a property of the type of situation, capturable in a rule that applies to all situations of that type.

\subsection{The state-dependence that escapes specification}
\label{p22:sub:eng-state-dependence}

Field-building is not like this, and the reason is the state-dependence of appropriateness established in the section on the coupling mechanism. The appropriateness of a gesture, a rhythm, an arrangement is not a property of the type of situation but of its effect on the holonomy of this field in this state, and the same gesture that builds the field in one state flattens it in another. There is no specification of the right placement of a thing, because the right placement depends on the state of the coupling it enters, and the state of the coupling is not a parameter that can be read off and fed into a procedure but a condition that must be perceived, freshly, in the concrete. What would be appropriate this evening, given how the day has gone, given the particular weather between the two, is not derivable from a rule about evenings, because the rule about evenings cannot contain the state of this evening's field, and the state is precisely what determines the appropriate. Engineering requires that the appropriate be a property of the type; field-building confronts an appropriate that is a property of the token, the particular field in its particular state, and a property of the token is exactly what no specification over types can reach.

\claim{Field-building escapes engineering because the appropriateness it must achieve is state-dependent and cannot be standardised. Engineering solves problems by specifications that apply to types of situation; the appropriateness of a field-building act is a property not of the type but of the act's effect on the holonomy of this field in this state, and the same act builds the field in one state and flattens it in another. No specification over types can capture a property of tokens, and so no engineering procedure can determine the appropriate placement, rhythm, or arrangement. What determines it can only be perceived in the concrete, by a sensibility resolved upon the state of the field, which is to say that field-building requires the very perceptual discrimination that distinguishes the aesthetic.}

This is the argument's core, and it can be stated as a general principle that the rest of the section will defend and qualify.

\begin{casebox}[The dinner that an optimiser would get wrong]
Imagine an engineer of the shared evening, equipped with the best specification a study of evenings could yield: the optimal time to eat, the arrangement of the table that maximises ease, the rhythm of conversation that the data show to be most satisfying, the lighting calibrated to the mood that surveys report as ideal. Apply the specification faithfully, evening after evening, and it will sometimes build the field and sometimes flatten it, and the specification cannot tell which, because what it cannot contain is the state of this evening. Tonight one of the two needs the meal early and quiet, the day having emptied her, and the specification's ideal rhythm of lively conversation is exactly wrong; another night the same liveliness is exactly right, the day having left them both with more to say than the evening can hold. The optimal lighting deflates an evening that wanted dimness for a grief being carried and lifts one that wanted brightness for a thing to celebrate. No refinement of the specification reaches this, because the specification is over types of evening and the appropriateness is a property of this evening's state, which the partners must perceive and the optimiser cannot encode. The engineer of the evening fails not because his specification is too coarse but because the thing he is trying to specify is not the kind of thing a specification can reach, a property of the token and not the type.
\end{casebox}

Wherever appropriateness cannot be exhausted by rules and can only be discerned by a trained perception of the concrete, the domain is aesthetic in the broad sense the genealogy recovered, the sense in which taste was a discernment of the fitting where no rule applies. Field-building is such a domain, because its appropriateness is state-dependent, and so field-building is aesthetic in this broad sense, not because it produces pretty results but because the discernment it requires is the discernment of the fitting where no rule decides, which is what aesthetic perception, at its root, has always been. The move from the field to aesthetics is therefore not an analogy or an ornament; it is the recognition that field-building has the formal mark of the aesthetic, the irreducibility of its appropriateness to rule.

\subsection{The strongest objection: architecture and design}
\label{p22:sub:eng-architecture}

The strongest objection points to disciplines that seem to be engineering and aesthetic at once. Architecture and design build conditional environments, arrange space and rhythm and atmosphere, and do so with a body of transmissible knowledge, patterns and principles and standards, that looks like engineering, while also exercising and demanding aesthetic judgement. If these disciplines combine engineering and aesthetics in the building of environments, then perhaps field-building is likewise a hybrid, partly specifiable and partly discerned, and the claim that it is aesthetic rather than engineering is overstated. The objection deserves to be met at its strongest, and meeting it requires looking at what the best reflection within these disciplines has itself said about the limits of specification, because the answer is not that architecture is merely engineering after all, but that architecture, rightly understood, has confronted the same irreducibility this paper describes and has named it.

The most searching account of design within the tradition arrived precisely at a quality that its own patterns could not specify, a quality without a name that the patterns could point toward and prepare for but never guarantee~\parencite{alexander_timeless}, present in the buildings and places that are alive and absent from those that are dead, and recognisable only to a perception that the patterns could train but not replace. The reason the quality is nameless repays attention, because it is the same reason field-building escapes engineering. The quality is nameless not because no word has yet been found for it but because it is not a property of the kind that words for properties name; it is not symmetry or proportion or any feature a building has in itself, since a building can have all the nameable virtues and be dead, and another can lack them and be alive. What the quality tracks is whether the place is alive, whether it accumulates rather than merely sits, whether being in it gains or flattens, and this is a feature of the relation between the place and those who live in it, in the same way that appropriateness is a feature of the relation between a gesture and the field it enters. The quality is nameless because it is relational and state-dependent, because it cannot be read off the building any more than appropriateness can be read off the gesture, and the deepest design thought reached this and named the namelessness rather than pretending to a specification it could not have.

The patterns, on this account, are not specifications that produce the quality; they are aids to a perception that must still discern, in the concrete case, whether the quality is present, and they work by cultivating that perception rather than by substituting for it. This is not the refutation of the present claim but its confirmation from within the discipline that seemed to threaten it. The deepest design thought reached the same conclusion the present section argues, that the appropriateness of a built environment cannot be specified and can only be discerned, that the transmissible knowledge of patterns serves to train a perception rather than to replace it, and that the quality the building must have is nameless precisely because it is not a property capturable in a rule. Architecture does not show that field-building is engineering. It shows that the building of living environments, pursued to its depth, becomes the cultivation of a perception of the unspecifiable fitting, which is what this paper says field-building is. The convergence is striking and worth marking: a tradition that set out to specify the building of good places, equipped with every tool of design, arrived by its own most rigorous path at the namelessness of the quality that matters, which is the same namelessness this paper finds in appropriateness, reached from the other side. Two disciplines, design and the philosophy of intimacy, pressing toward the conditions of a living environment, meet at the same recognition, that the quality of a place or a field is relational, state-dependent, and discernible only by a trained perception, never specifiable in advance.

\subsection{The transmissible and the discerned}
\label{p22:sub:eng-transmissible}

The architecture objection, met, yields a refinement the account needs, concerning what can and cannot be transmitted in the cultivation of field-building. It is not that nothing about field-building can be taught, that there are no patterns, no transmissible knowledge, no aids to perception; there are, and the four-step model and its practices are themselves such aids. The claim is narrower and exact: what can be transmitted are aids to a perception that must still discern the concrete case, and what cannot be transmitted is the discernment itself, the perceiving of whether this field in this state is being built or flattened, which no aid can perform on the perceiver's behalf. This is the structure of every aesthetic cultivation, and it is why such cultivation is an education and not an instruction, a raising of a perceptual power by exercise rather than a transmission of a code by telling. The patterns prepare the perception; the perception does the discerning; and the gap between the prepared and the discerned is the gap between engineering and aesthetics, the gap that the state-dependence of appropriateness opens and that no quantity of transmissible knowledge can close. Field-building, then, has exactly the structure the rest of the paper requires: it is cultivable, since perception can be trained, and it is not specifiable, since the trained perception must still discern the token, and a capacity that is cultivable but not specifiable is the proper object of an aesthetic education, which is the name for the cultivation of a discerning perception of the fitting where no rule decides.
